\section{Capítulo~\ref{sec:livingmachine}. A máquina de neurônios vivos}

O comportamento das culturas em matrizes de eletrodos foi medido em experimentos diretos com registro e estimulação; o mecanismo das salvas se apoia no modelo de esgotamento das sinapses e em experimentos com microculturas, e as afirmações sobre a dopamina na placa, em experimentos isolados, parte dos quais os próprios autores consideram apenas uma demonstração.

{\footnotesize
\setlength{\tabcolsep}{3pt}\renewcommand{\arraystretch}{1.15}
\begin{longtable}{>{\raggedright}p{0.5cm}>{\raggedright}p{4.0cm}>{\raggedright}p{5.6cm}>{\raggedright}p{2.3cm}>{\raggedright\arraybackslash}p{1.7cm}}
\toprule
N.º & Proposição & Experimento e o que foi medido & Fonte & Status \\
\midrule\endfirsthead
\toprule
N.º & Proposição & Experimento e o que foi medido & Fonte & Status \\
\midrule\endhead
1 & Cultura no chip: sobretudo \nt{GLUT}, com uma fração menor de \nt{GABA} que depende da preparação; em culturas de hipocampo, cerca de $6\%$ & Uma rede de milhares de neurônios sobre uma matriz de eletrodos é um experimento padrão (linha~3). A fração de neurônios \nt{GABA} em culturas de hipocampo foi contada por marcação para GABA: cerca de $6\%$ dos neurônios. Para culturas de córtex não damos um número verificado & \cite{Benson1994} & medido \\
2 & Organoides: tecido nervoso tridimensional com cerca de um milímetro, a partir de células-tronco humanas & A partir de células-tronco foram cultivados organoides tridimensionais com várias regiões cerebrais, entre elas um córtex com camadas de precursores e neurônios maduros & \cite{Lancaster2013} & medido \\
3 & Sem entrada, a cultura dispara em salvas com intervalos de um segundo a minutos, conforme a cultura & $58$ culturas de córtex com densidade de $3\,000$ a $50\,000$ neurônios foram registradas por cinco semanas ($963$ registros de meia hora): depois de duas semanas, em quase todas as culturas predominam salvas da rede inteira; os intervalos entre salvas vão de $1$ a $300$~s e dependem fortemente da preparação & \cite{Wagenaar2006} & medido \\
4 & A salva termina por esgotamento do neurotransmissor, intervalo $T_{\text{salva}}\approx\tau_{\text{rec}}\ln\frac{1}{1-x^*}$~\eqref{eq:burstinterval}; $2$--$4$~s é a estimativa do modelo com $\tau_{\text{rec}}=0.8$~s, a recuperação medida é mais lenta & A fórmula é deduzida no capítulo. Modelo: uma rede com sinapses que se esgotam gera sozinha salvas da rede inteira. Experimento em microculturas de $4$--$30$ neurônios: a duração da salva depende do tempo desde a anterior, e o esgotamento das vesículas de neurotransmissor elimina as salvas; conclusão dos autores: a salva é limitada pelo estoque de neurotransmissor. Lá a recuperação do estoque leva dezenas de segundos, o que, com a mesma fórmula, dá intervalos de dezenas de segundos e minutos & dedução no capítulo; \cite{Tsodyks2000,CohenSegal2011,Tsodyks1997} & teoria \\
5 & ``O professor que se cala'': a estimulação é desligada quando aparece a resposta desejada, e a resposta se fixa & A cultura foi estimulada a $0.3$--$1$~Hz até aparecer a resposta escolhida $50\pm10$~ms depois do estímulo; então a estimulação era desligada por $5$~min; depois de alguns ciclos a resposta desejada surgia de imediato; foram traçadas curvas de aprendizado & \cite{Shahaf2001} & medido \\
6 & A cultura joga tênis; aumento significativo das sequências na sessão só com realimentação completa & Em detalhe no capítulo~\ref{sec:dishbrain} e em seu suplemento: com silêncio no lugar do estímulo, a cultura também joga melhor no fim da sessão do que sem realimentação, mas com efeito menor; o aprendizado de uma sessão para outra é instável & \cite{Kagan2022} & medido \\
7 & Que a cultura aprende a prever sua entrada é hipótese; as palavras fortes sobre ``inteligência'' foram contestadas & O princípio da energia livre é uma proposta teórica; não há experimento direto que o distinga de explicações mais simples. Trinta pesquisadores, num artigo de resposta, apontaram que a mudança de comportamento na malha não mostra nem inteligência nem sensações & \cite{Friston2010,Balci2023} & hipótese \\
8 & A cultura conduz um corpo virtual até o alvo com um padrão de estímulos escolhido & O programa escolhia sozinho os padrões de estimulação de treino pelo resultado corrente; em dezenas de minutos a rede atingia os estados pedidos, e a plasticidade durava $80$~min após $2$~h de treino. Os mesmos estímulos, aplicados sem relação com o comportamento, não tinham efeito & \cite{Bakkum2008} & medido \\
9 & Reservatório: só se treina a leitura linear $y=W_{\text{saída}}\mathbf r$~\eqref{eq:reservoir} & Teoria da computação de reservatório: qualquer sistema não linear com memória que se apaga, mais uma saída linear treinada & \cite{Maass2002,JaegerHaas2004} & teoria \\
10 & O organoide como reservatório distingue falantes com acurácia de cerca de $78\%$ (segundo os autores); a leitura aprende pelo erro, e as conexões do organoide mudam sem professor & Sons de fala de vários falantes foram aplicados ao organoide como padrões de corrente na matriz de eletrodos; a leitura treinada distinguia o falante com acurácia de cerca de $78\%$, segundo os autores. Eles também relatam que durante o treino a conectividade funcional do organoide mudou, e chamam isso de aprendizado não supervisionado no próprio organoide & \cite{Cai2023} & medido \\
11 & Um milhão de neurônios gasta cerca de $10^{-4}$~W em disparos~\eqref{eq:dishpower} & Estimativa: o custo de um disparo, $10^{-10}$~J, do orçamento energético do córtex~\eqref{eq:energy}, multiplicado pelo número de disparos; não há medição direta da potência da cultura & dedução no capítulo; \cite{Attwell2001} & teoria \\
12 & Dopamina por clarão de luz: $1$~mM de dopamina enjaulada, $240$ repetições, o número de disparos evocados cresceu & Numa plataforma remota de organoides, dopamina numa gaiola fotossensível ($1$~mM) era liberada com ultravioleta ($365$~nm, $0.8$~s) quando o estímulo evocava mais disparos; em $240$ passos (cerca de $5$~h) o número de disparos em geral cresceu. Os autores escrevem que um único experimento não permite estabelecer a relação com a dopamina; não há controles & \cite{Jordan2024} & hipótese \\
13 & A dopamina de células vivas muda o peso de um transistor orgânico de forma duradoura e reversível & Células que liberam dopamina foram cultivadas sobre um elemento neuromórfico orgânico; a oxidação da dopamina junto ao eletrodo muda a condutância do canal; foram mostrados a mudança duradoura do peso e seu retorno & \cite{Keene2020} & medido \\
14 & Existem organoides com neurônios \nt{DOPA} funcionais; sua dopamina não foi usada como professor & Em organoides de células-tronco humanas foram encontrados neurônios \nt{DOPA} maduros e eletricamente ativos, e produção de dopamina. Não encontramos na literatura experimentos em que essa dopamina ensine outra rede & \cite{Jo2016} & medido \\
15 & O organoide equilibra a haste: os estímulos de ensino do programa são melhores que os aleatórios, sem consolidação, por meio das sinapses \nt{GLUT} & Organoides de córtex de camundongo em malha fechada com a tarefa do ``pêndulo no carrinho'': na maioria dos organoides, os sinais escolhidos por aprendizado por reforço deram resultado melhor que os aleatórios ou nenhum; após $45$~min de repouso a melhora não se mantinha; o bloqueio dos receptores AMPA e NMDA a anulava & \cite{Robbins2026} & medido \\
16 & Sem registradores de controle, a rede na bancada não consegue decidir sozinha o que conta como sucesso (proposição~\ref{prop:livingmachine}) & Em todos os experimentos das linhas 5--15, o sinal de ensino é definido pelo experimentador ou por um programa; não há experimento em que a cultura tenha criado sozinha um critério de sucesso; é uma conclusão pela ausência, não um experimento & --- & hipótese \\
\bottomrule
\end{longtable}}
